\SourceLeaf{341}

\begin{center}
[3.]
\end{center}

Schliesslich wollen wir zeigen, wie man von dem Beweis in [1] auf den in [2] kommen kann, indem
wirklich der Beweis zuerst durch Laterale gefunden ist, und erst später durch die geometrische
Betrachtung.

Untersuchen wir, wie sich der Beweis in [1] stellt, wenn ein Punkt des Vierecks, z.\,B. (4), mit
dem Anfangspunkt der Zählung coincidirt. Man hat dann $p^{\mathrm{IV}} = 0$, also sind die
Coordinaten
\[
\begin{array}{rl}
\text{von I:} & p'p'' \\
\text{von II:} & p'p''' \\
\text{von III:} & p''p'''.
\end{array}
\]
Also sind nach [2] die Seiten:
\[
\begin{array}{rl}
(\text{I}, \text{II}) & = (4,1) . (4,2) \\
(\text{II}, \text{III}) & = (4,1) . (4,3) \\
(\text{III}, \text{I}) & = (4,2) . (4,3).
\end{array}
\]
Folglich verhalten sich die Seiten des Dreiecks
\[
(\text{I}, \text{II}) : (\text{II}, \text{III}) : (\text{III}, \text{I}) = (4,1) . (4,2) :
(4,1) . (4,3) : (4,2) . (4,3),
\]
wie wir sie in [2] durch Construction bestimmten.
